\chapter{진짜 공책을 펼치다}

지금까지는 교실 다섯 명으로 연습했어요.
이제 진짜 블록체인 공책을 펼칠 차례예요.

\section{아주 큰 도서관}

제가 쓴 공책은 \textbf{Arbitrum One}이라는 블록체인이에요.
이 공책은 너무 커서 집 컴퓨터로는 다 펼쳐 볼 수 없어요.
다행히 구글이 이 공책의 사본을 \textbf{BigQuery}라는 곳에 정리해 두었어요.
BigQuery는 아주 큰 도서관 같은 곳이에요.
사서에게 “이런 줄만 찾아 주세요”라고 적어서 내면 사서가 찾아 줘요.
이 쪽지를 \textbf{질의}라고 해요.

그런데 이 도서관은 공짜가 아니에요.
사서가 책장을 많이 넘길수록 돈을 내야 해요.
그래서 저는 질의 하나가 넘길 수 있는 양에 상한을 두었어요.
질의마다 먼저 “얼마나 넘겨야 하나요?”라고 물어보고(연습 질의),
150기가바이트가 넘으면 그 질의는 하지 않도록 프로그램을 만들었어요.

\section{누구의 기록을 볼까?}

Arbitrum One에는 지갑이 수백만 개 있어요.
그 모두를 볼 수는 없어요.
그래서 먼저 “누구를 볼지”를 정했어요.

제가 고른 사람들은 \textbf{GMX}라는 앱을 쓰는 사람들이에요.
GMX는 어떤 가격이 오를지 내릴지 맞히는 거래를, 빌린 돈까지 보태서 크게 할 수 있는 앱이에요.
잘 맞히면 많이 벌지만, 틀리면 많이 잃어요.
너무 많이 잃어서 맡겨 둔 돈(담보)이 모자라면 앱이 거래를 강제로 끝내 버리는데, 이것을 \textbf{청산}이라고 해요.
어른들에게도 아주 위험한 일이에요.
이 책은 이런 거래를 권하지 않아요.
저는 이 사람들이 남긴 \emph{기록}만 연구에 썼어요.

왜 GMX를 쓰는 사람들을 골랐을까요?
GMX 기록에는 허락 지도에도 송금 지도에도 없는 정보가 있기 때문이에요.
거래를 몇 번 닫았는지, 그중 몇 번이 청산이었는지 같은 것들이에요.
이런 정보는 나중에 점수를 채점할 때 두 지도와 상관없는 “제3의 기록”이 되어 줘요.

\section{연구 기간과 대상}

연구 기간은 2025년 12월 1일부터 2026년 5월 31일까지 여섯 달이에요.
이 기간에 GMX에서 거래를 한 번이라도 닫은 지갑은 9,388개였어요.
그중 거래를 \textbf{세 번 이상} 닫은 지갑 \textbf{5,521개}를 골랐어요.
이 지갑들을 \textbf{매칭 코호트}라고 불러요.
코호트는 “같은 조건으로 함께 묶은 무리”라는 뜻이에요.

그다음, 이 5,521개 지갑이 허락을 하거나 받은 줄, 송금을 하거나 받은 줄을
한 달씩 나누어 도서관에서 찾아왔어요.

\begin{center}
\small
\begin{tabular}{@{}lr@{}}
\toprule
가져온 기록 (2025년 12월 \textasciitilde{} 2026년 5월) & 줄 수 \\
\midrule
허락 줄 & 453,545 \\
송금 줄 & 2,801,988 \\
GMX 거래 닫기 줄 & 365,488 \\
\bottomrule
\end{tabular}
\end{center}

\section{달력으로 보는 이 연구}

이 책의 뒷부분은 날짜가 아주 중요해요.
그래서 달력을 하나 그려 둘게요.
지금은 다 이해하지 못해도 괜찮아요.
13장과 14장을 읽을 때 이 쪽으로 돌아와 보세요.

\begin{figure}[h]
\centering
\begin{tikzpicture}[x=0.92cm, y=1cm]
  % 달 칸
  \foreach \i/\m in {0/12월, 1/1월, 2/2월, 3/3월, 4/4월, 5/5월, 6/6월, 7/7월, 8/8월} {
    \draw[greycol!60] (\i,0) rectangle (\i+1,0.55);
    \node[font=\scriptsize] at (\i+0.5,0.275) {\m};
  }
  \node[font=\scriptsize, anchor=east] at (-0.1,0.275) {2025};
  \node[font=\scriptsize] at (1.5,0.85) {2026};
  % 연구 기간
  \draw[decorate, decoration={brace, amplitude=4pt}, line width=0.8pt] (0,1.15) -- node[above=4pt, font=\scriptsize] {연구 기간 (6개월)} (6,1.15);
  % 봄 홀드아웃
  \fill[allowcol!25] (0,-0.75) rectangle (3,-0.35);
  \fill[allowcol!60] (3,-0.75) rectangle (6,-0.35);
  \node[font=\scriptsize] at (1.5,-0.55) {점수 매기기};
  \node[font=\scriptsize, color=white] at (4.5,-0.55) {채점 (봄)};
  \node[font=\scriptsize, anchor=east] at (-0.1,-0.55) {13장};
  % 등록 재현
  \fill[sendcol!25] (0,-1.45) rectangle (6,-1.05);
  \fill[sendcol!75] (6,-1.45) rectangle (9,-1.05);
  \node[font=\scriptsize] at (3,-1.25) {점수 매기기};
  \node[font=\scriptsize, color=white] at (7.5,-1.25) {채점 (여름)};
  \node[font=\scriptsize, anchor=east] at (-0.1,-1.25) {14장};
  % 동결선
  \draw[dashed, line width=0.8pt] (3,0.6) -- (3,-1.7) node[below, font=\scriptsize, align=center] {2월 28일\\봉투 봉함};
  \draw[dashed, line width=0.8pt] (6,0.6) -- (6,-1.7) node[below, font=\scriptsize, align=center] {5월 31일\\봉투 봉함};
\end{tikzpicture}
\caption{연구 기간과 두 번의 “미래 맞히기” 시험. 점선은 점수를 봉투에 넣고 봉한 날이에요.}
\end{figure}

\begin{deeper}[진짜 이름과 숫자]
\begin{itemize}
\item BigQuery의 공개 테이블 \filepath{bigquery-public-data.goog_blockchain_arbitrum_one_us.logs}에서 기록을 가져왔어요.
\item GMX V2의 \texttt{PositionDecrease} 이벤트(포지션을 줄이거나 닫는 기록, 청산 포함)를 “거래 닫기”로 셌어요.
      거래한 사람은 거래를 실행한 프로그램의 주소가 아니라 이벤트에 적힌 계정(\texttt{account})으로 알아봤어요.
\item 질의 하나의 상한은 150 GiB였고, 질의마다 연습 질의(dry run)로 먼저 크기를 확인했어요.
\item 송금 줄 가운데 보낸 양이 0인 줄을 빼고 남은 2,639,388줄로 같은 기간의 송금 지도를 그렸어요.
\item 같은 지갑 목록으로 Aave라는 대출 앱의 기록도 가져왔지만, 이 기록은 논문 본문의 결과에는 쓰지 않았어요.
\end{itemize}
\end{deeper}

\begin{learned}
\begin{itemize}
\item 진짜 기록은 Arbitrum One 블록체인에서, BigQuery라는 큰 도서관을 통해 가져왔어요.
\item 2025년 12월부터 2026년 5월까지 GMX에서 거래를 세 번 이상 닫은 지갑 5,521개를 연구했어요.
\item 이 지갑들의 허락 줄 약 45만 개와 송금 줄 약 280만 개로 두 지도를 그렸어요.
\end{itemize}
\end{learned}
